\documentclass[12pt]{article}
\usepackage{amsmath,amssymb}

\begin{document}
\section*{{2024 AMC 10A Problems/Problem 19}}
Source: AoPS Wiki

\subsection*{Solution}
For a geometric sequence, we have $ab=720^2=2^8 3^4 5^2$ , and we can test values for $b$ . We find that $b=768$ and $a=675$ works, and we can test multiples of $5$ in between the two values. Finding that none of the multiples of 5 divide $720^2$ besides $720$ itself, we know that the answer is $7+6+8=\boxed{\textbf{(E) } 21}$ . (Note: To find the value of $b$ without bashing, we can observe that $2^8=256$ , and that multiplying it by $3$ gives us $768$ , which is really close to $720$ . ~ YTH) Note: The reason why $ab=720^2$ is because $b/720 = 720/a$ . Rearranging this gives $ab = 720^2$ ~eevee9406 Note: Another reason that $ab=720^2$ is because the $\sqrt{ab}=720$ (as the middle term in a geometric series is always the geometric mean [the geometric mean is the square root of the product of the first and last terms of the series]) and squaring on both sides results in $ab=720^2$ . ~ThatPrimePunnyGuy Note: Because it is a geometric sequence, it is clear that $a*k = 720$ and $720*k = b$ . Then, $a*k = 720 = b/k$ . Multiply the first and third term together and it is equal to the middle term multiplied by itself. ~RON-WEASLEY

\end{document}
